\subsection{Triangularisable endomorphisms}

In this section we will be interested in the case where the minimal polynomial \(p(\mathbf{X})\) of u splits into a product of linear factors in K{[}X{]}, in other words (VII, p.~33, Cor. 2) in the case where all the eigenvalues of u belong to K. This will hold in particular when K is algebraically closed. Prop. 3 of VII, p.~31 gives immediately :

PROPOSITION 7. --- Let E be a finite dimensional vector space over a commutative field K, and let u be an endomorphism of E whose eigenvalues all belong to K. For each eigenvalue \(\alpha\) of u, let \(M_{\alpha}\) be the vector subspace of E consisting of those elements x for which there exists an integer \(k \geqslant 1\) such that \((u - \alpha)^{k} \cdot x = 0\) . Then \(M_{\alpha}\) is closed under u, the vector space E is the sum of the \(M_{\alpha}\) , and there exist polynomials \(s_{\alpha} \in K[X]\) such that, for all \(x \in E\) , the component of x in \(M_{\alpha}\) is equal to \(s_{\alpha}(u) \cdot x\) .

The submodule \(M_{\alpha}\) , being finitely generated as a K{[}X{]}-module, then has an annihilator of the form \((\mathbf{X}-\alpha)^{r}\) ; in other words, there exists an integer \(r \geqslant 1\) such that

\[
(u - \alpha) ^ {r} \cdot x = 0
\]

for all \(x \in \mathbf{M}_{\alpha}\) ; the restriction of \(u - \alpha\) to \(\mathbf{M}_{\alpha}\) is a nilpotent endomorphism.

Still assuming that the eigenvalues of u belong to K, we now apply Prop. 6 of VII, p.~34 to u. The polynomials \(p_{k}\) are nothing other than the X - \(\alpha\) (as \(\alpha\) runs through the set of eigenvalues of u), and we see that E is the direct sum of subspaces \(E_{i}\) closed under u, cyclic (with respect to u), and such that the minimal polynomial of the restriction of \(u\) to \(E_i\) has the form \((X - \alpha)^m\) . Let \(E_i'\) be the K{[}X{]}-module associated to \(E_i\) ; then \(E_i'\) is isomorphic to one of the modules K{[}X{]}/\((X - \alpha)^{m}\). Now the residue classes mod(X - \(\alpha\) ) \(^m\) of the elements (X - \(\alpha\) ) \(^k\) ( \(0 \leqslant k \leqslant m - 1\) ) form a K-basis of K{[}X{]}/((X - \(\alpha\) ) \(^m\) ) (IV, p.~11, Cor.), and

\[
\mathbf {X} (\mathbf {X} - \alpha) ^ {k} = \alpha (\mathbf {X} - \alpha) ^ {k} + (\mathbf {X} - \alpha) ^ {k + 1}
\]

for \(0 \leqslant k \leqslant m - 1\) ; it follows that \(\mathbf{E}_i\) has dimension \(m\) , and if \(\alpha\) is the unique eigenvalue of the restriction \(u_i\) of \(u\) to \(\mathbf{E}_i\) , then there exists a basis of \(\mathbf{E}_i\) with respect to which the matrix of \(u_i\) is the \(m \times m\) matrix

\[
U _ {m, \alpha} = \left( \begin{array}{c c c c c c} \alpha & 0 & 0 & ... & 0 & 0 \\ 1 & \alpha & 0 & ... & 0 & 0 \\ 0 & 1 & \alpha & ... & 0 & 0 \\ \cdots & \cdots & \cdots & \cdots & \cdots & \cdots \\ 0 & 0 & 0 & ... & \alpha & 0 \\ 0 & 0 & 0 & ... & 1 & \alpha \end{array} \right)\tag{4}
\]

DEFINITION 3. --- For every field K, every integer \(m \geqslant 1\) , and every \(\alpha \in K\) , the matrix \(U_{m,\alpha}\) is called the Jordan matrix of order m and eigenvalue \(\alpha\) .

PROPOSITION 8. --- Let E be a finite dimensional vector space over a commutative field K, and let u be an endomorphism of E. Then the following conditions are equivalent :

\begin{enumerate}
\def\labelenumi{(\roman{enumi})}
\item
  the eigenvalues of \(u\) (in some algebraically closed extension of \(\mathbf{K}\) ) belong to \(\mathbf{K}\) ;
\item
  there exists a basis of \(\mathbf{E}\) with respect to which the matrix of \(u\) is lower (resp. upper) triangular;
\item
  there exists a basis of E with respect to which the matrix of u is a diagonal block of Jordan matrices.
\end{enumerate}

We have (i) \(\Rightarrow\) (iii) by Prop. 7 and the above remarks, and the assertions (iii) \(\Rightarrow\) (ii) and (ii) \(\Rightarrow\) (i) are trivial.

DEFINITION 4. --- An endomorphism satisfying conditions (i), (ii) and (iii) of Prop. 8 is called triangularisable.

In particular if K is algebraically closed, then every endomorphism of a K-vector space is triangularisable.

For matrices, Prop. 8 implies :

COROLLARY. --- Let U be a square matrix over a commutative field K such that all the eigenvalues of U are in K; then there exists a matrix similar to U which is a diagonal block of Jordan matrices.

Remarks. --- 1) It follows from Prop. 6 of VII, p.~34, that, if U is similar to a diagonal block of Jordan matrices \((J_{k})\) , then the number of \(J_{k}\) of the form \(U_{m,\alpha}\) (for given m and \(\alpha\) ) is uniquely determined by U.

\begin{enumerate}
\def\labelenumi{\arabic{enumi})}
\setcounter{enumi}{1}
\tightlist
\item
  More generally, if \(U\) is similar to a diagonal block of Jordan matrices \(U_{m_i,\alpha_i}\) , then the similarity invariants of \(U\) can be readily calculated by a method modelled on that presented in VII, p.~25, Remark 3: write all the \((X - \alpha_i)^{m_i}\) with the same \(\alpha\) on the same line, in decreasing order of exponents, and complete with 1's to have lines whose lengths are equal to the order of \(U\) ; this done, the similarity invariants of \(U\) are obtained, in decreasing order of indices, by forming the products of terms in the same column. For example, for the matrix
\end{enumerate}

\[
\left( \begin{array}{c c c} 2 & 0 & 0 \\ 0 & 3 & 0 \\ 0 & 1 & 3 \end{array} \right)
\]

we write

\[
\begin{array}{l} (\mathbf {X} - 2), 1, 1 \\ (\mathbf {X} - 3) ^ {2}, 1, 1 \end{array}
\]

and the similarity invariants are 1, 1 and \((\mathbf{X} - 2)(\mathbf{X} - 3)^2\) .

By noticing that the minimal polynomial of the Jordan matrix \(U_{m,\alpha}\) is \((X-\alpha)^{m}\) , and that it is equal to its characteristic polynomial, we obtain the following result:

PROPOSITION 9. --- If the square matrix U is similar to a diagonal block of Jordan matrices \((U_{m_{i},\alpha_{i}})\) , then the minimal polynomial of U is the lcm of the \((\mathbf{X}-\boldsymbol{\alpha}_{i})^{m_{i}}\) , and the characteristic polynomial is the product of the \((\mathbf{X}-\boldsymbol{\alpha}_{i})^{m_{i}}\) .

COROLLARY. --- In the notation of Prop. 7, the dimension of the subspace \(M_{\alpha}\) is the multiplicity of the eigenvalue \(\alpha\) as a root of the characteristic polynomial of u.

